% =============================================================================
% Section 03 — Expression and Determination
% Book I — Ontology of Viability
% Chapter 04 — Expression
%
% Purpose:
%   To examine the ontological relation between expression
%   and determination without reducing either concept to the other.
%
% Conceptual Policy:
%   1. Expression and determination are analytically distinct.
%   2. Not every determination is necessarily an expression.
%   3. Not every expression is necessarily a determination.
%   4. Expression must not be reduced to the articulation or
%      disclosure of determinations.
%   5. Determination must not be defined through expression alone.
%   6. Their relation must be examined according to subject
%      and respect.
% =============================================================================

\section{Expression and Determination}
\label{sec:expression-and-determination}

\subsection{The Ontological Problem of Expression and Determination}

The relation between expression and determination constitutes
a further problem in the ontological analysis of expression.

An existent may possess determinations through which
its ontological constitution is specified, delimited,
or differentiated. It may also be associated with
expressions that raise a distinct question concerning
how aspects of its being are articulated.

The presence of both determination and expression
does not establish that they are identical or that
one can be reduced to the other.

Determination concerns the ontological specification
of an existent in a particular respect. Expression
concerns the ontological status of an articulation
or expression associated with an existent. The precise
meaning of expression remains under examination,
but its distinction from determination must be
preserved throughout that inquiry.

The central problem is therefore to establish whether
expression is a determination, depends upon determinations,
articulates determinations, or stands in another
ontological relation to them.

These possibilities must not be treated as equivalent.
Each implies a different account of the constitution
of an existent and of the role that expression may
occupy within it.

\subsection{Determination as Ontological Specification}

Determination concerns the specification of an existent
in a particular ontological respect.

An existent is not adequately described merely by
asserting that it exists. Its determinations may
establish particular ways in which it is constituted,
distinguished, related, or delimited.

Determination therefore belongs to the analysis
of what specifies an existent, rather than necessarily
to the analysis of how that existent is expressed.

A determination may be ontologically significant
even when no corresponding expression has been
identified. The absence of an identified expression
does not establish the absence of determination.

Likewise, the identification of a determination
does not automatically establish that the determination
has been expressed.

This distinction prevents determination from being
made dependent upon the occurrence, recognition,
or description of expression.

Determination must therefore be examined according
to its own ontological conditions, even where it
is subsequently found to be related to expression.

\subsection{Expression Cannot Be Presumed to Be a Determination}

One possible interpretation treats expression as
a determination of an existent.

According to this interpretation, to express would
be one of the ways in which an existent is determined.
Expression would consequently occupy a position
within the ontological constitution of the existent
alongside other determinations.

This possibility cannot be excluded in advance.
However, it cannot be accepted merely because
expression is associated with an existent.

The fact that something belongs to, proceeds from,
or is related to an existent does not establish
that it is a determination of that existent.

To classify expression as a determination would
require an account of what expression specifies,
delimits, or establishes in the existent and
why this function should be understood as
determination rather than through another
ontological category.

Furthermore, if expression is treated as a determination,
it becomes necessary to distinguish between the
determination that makes expression possible and
the expression itself.

Otherwise, the concept risks becoming circular:
expression would be explained as a determination
while the relevant determination would itself
be identified only through expression.

The classification of expression as a determination
must therefore remain a possible interpretation,
not an established premise.

\subsection{Not Every Determination Is Necessarily Expressed}

The relation between determination and expression
must also be examined from the opposite direction.

Even if certain expressions are found to articulate
or correspond to determinations, it does not follow
that every determination must be expressed.

A determination may belong to the ontological
constitution of an existent without being associated
with an identifiable expression.

The possibility of unexpressed determination
must therefore remain conceptually available.

To deny this possibility would require establishing
that expression is a necessary condition for the
existence or ontological efficacy of every
determination.

Such a requirement has not yet been demonstrated.

Moreover, the distinction between the existence
of a determination and its expression prevents
the ontological status of the former from being
made dependent upon the latter.

A determination does not cease to be ontologically
relevant merely because no corresponding expression
has been established.

This does not imply that determination is always
independent of expression. It means only that
their dependence cannot be presumed to be universal.

\subsection{Not Every Expression Is Necessarily a Determination}

The converse restriction is equally important.

Even where an expression is established, it does
not follow that the expression itself constitutes
a determination of the existent.

An expression may be related to a determination
without being identical to it. It may articulate
a determination, arise under conditions established
by determinations, or stand in another relation
that does not make it a determination in its own right.

The distinction is necessary because the classification
of every expression as a determination would eliminate
the possibility that expression has a distinct
ontological status.

Such a reduction would also make the analysis
of expression dependent upon the conceptual
framework of determination, leaving no independent
question concerning what expression itself is.

An expression may therefore be examined as an
expression even when its relation to determination
has not been established.

Whether a particular expression is itself a
determination must be decided according to
the ontological role it occupies, not according
to the mere fact that it is associated with
an existent.

\subsection{Expression as Articulation of Determination}

A further possibility is that expression articulates
or gives form to certain determinations of an existent.

Under this interpretation, determination and
expression remain distinct, while expression
is understood through a relation to what has
already been ontologically specified.

This possibility may explain why an existent can
possess determinations that are not exhausted
by the particular expressions associated with it.

However, the term articulation must itself be
used carefully.

If articulation means that expression makes
a determination available to a recipient,
it introduces an epistemic or communicative
framework that does not belong to the present
ontological inquiry.

If articulation is intended in a strictly
ontological sense, its meaning must be specified
without presupposing a knower, interpreter,
observer, or receiving consciousness.

The ontological relation may concern the way
in which a determination is configured, related,
or expressed within the being of an existent.
Yet the exact nature of this relation remains
to be established.

Accordingly, expression as articulation of
determination is a possible account, but it
cannot serve as a complete definition of
expression without further argument.

\subsection{Expression and the Configuration of Determinations}

Determinations may occur in organized configurations.

An existent may possess several determinations
whose relations establish a particular structure
or configuration. These determinations may be
mutually dependent, hierarchically ordered,
or related through other forms of organization.

Expression may be associated with such configurations.
It may concern a particular determination, a relation
among determinations, or a configuration considered
as a whole.

These possibilities must remain distinct.

An expression associated with one determination
does not necessarily express the entire configuration.
Likewise, an expression associated with a configuration
does not establish that every determination within
that configuration is independently expressed.

The relation between expression and configuration
therefore requires specification of its subject
and respect.

A further distinction must be preserved between
the organization of determinations and the organization
of expressions. Even where these organizations
correspond, they are not necessarily identical.

The structure of an existent cannot be inferred
exhaustively from the structure of its expressions
without demonstrating the ontological adequacy
of that inference.

\subsection{Determination, Expression, and Ontological Dependence}

The relation between expression and determination
may involve different forms of ontological dependence.

Expression may depend upon certain determinations
because those determinations establish conditions
under which a particular expression can occur.

Alternatively, a determination may depend upon
expression if expression is shown to contribute
to the constitution or specification of the existent
in a particular respect.

A reciprocal relation may also be possible, although
such reciprocity cannot be assumed merely because
the two concepts are related.

The existence of a relation does not establish
its direction, symmetry, or degree of dependence.

It is therefore necessary to distinguish among:

\begin{enumerate}
    \item The dependence of expression upon determination.
    \item The dependence of determination upon expression.
    \item Reciprocal dependence between expression
    and determination.
    \item A relation between them that involves
    neither constitutive dependence nor equivalence.
\end{enumerate}

These alternatives have different ontological
implications.

In particular, the possibility that expression
depends upon determination does not establish
that expression is reducible to determination.
Dependence and identity are different relations.

Similarly, the possibility that expression contributes
to a determination does not establish that determination
is nothing beyond expression.

The precise form of dependence must therefore
be established independently in each relevant case.

\subsection{Expression and the Distinction Between Determination and Identity}

The relation between expression and determination
must also preserve the distinction between both
concepts and identity.

Determination does not exhaust identity, and
expression does not exhaust determination.

An existent may possess determinations through
which aspects of its constitution are specified.
It may also possess expressions related to some
of those determinations.

However, neither the totality of determinations
nor the totality of expressions can be assumed
to constitute the identity of the existent.

This restriction is particularly important because
the identification of an organized set of determinations
may create the appearance of a complete account
of an existent.

Yet a complete description of determinations
within a specified analytical domain is not
necessarily equivalent to an exhaustive account
of identity.

The same applies to expressions.

The ontological analysis must therefore preserve
three distinct questions:

\begin{enumerate}
    \item What determines or specifies the existent?
    \item What is expressed, and what is the ontological
    status of that expression?
    \item What constitutes the identity of the existent?
\end{enumerate}

The answers may be related, but they cannot be
treated as interchangeable.

\subsection{Expression, Determination, and Change}

The distinction between expression and determination
becomes particularly important when an existent
undergoes change.

An existent may undergo a change in one or more
of its determinations. It may also undergo a change
in its expressions without a corresponding change
in the determinations to which those expressions
are related.

Conversely, a determination may change while
a particular expression remains unchanged,
or while no corresponding expressive alteration
has been established.

These possibilities indicate that expressive change
and determinate change cannot be treated as
necessarily equivalent.

A change in expression may be associated with
a change in determination, but the association
does not establish identity between the two changes.

Likewise, the preservation of an expression does
not necessarily establish the preservation of
the determination associated with it.

The analysis must identify what has changed,
the respect in which the change occurred,
and the ontological relation between the relevant
expression and determination.

Without these distinctions, the concept of expression
would become an unreliable substitute for the
analysis of ontological change.

\subsection{The Limits of Inferring Determination from Expression}

An expression may provide grounds for considering
a possible determination, but the existence of
an expression does not automatically establish
the complete determination of the existent.

Such an inference would require an account of
the relation between the expression and the
determination attributed to the existent.

In particular, it would be necessary to establish
whether the expression is constitutive of the
determination, dependent upon it, or merely
associated with it.

The same restriction applies when several expressions
are considered together.

The accumulation of expressions does not necessarily
produce an exhaustive account of determinations.
Nor does an extensive configuration of expressions
establish that every relevant determination has
been accounted for.

The distinction between ontological determination
and the description of expression must therefore
be maintained.

This restriction does not deny that expressions
may be ontologically significant. It prevents
their significance from being extended beyond
what their established relations justify.

\subsection{Subject and Respect in the Relation Between Expression and Determination}

Every claim concerning the relation between
expression and determination must specify
its subject and respect.

A statement that an existent expresses a determination
is incomplete unless it clarifies which existent
is under consideration, which determination
is involved, and in what respect the expression
is related to it.

The same expression may be associated with
different determinations under different conditions.
Likewise, a single determination may be related
to different expressions across distinct domains.

Consequently, the relation cannot be generalized
without identifying the ontological scope within
which it is asserted.

The specification of subject and respect is
particularly important where multiple determinations
and expressions are involved.

Without it, distinct relations may be conflated,
and a local correspondence may be incorrectly
treated as a universal ontological principle.

The analysis of expression and determination
must therefore remain sensitive to the scope
and conditions of each proposed relation.

\subsection{Open Questions Concerning Expression and Determination}

The analysis of this section establishes a number
of distinctions, but it does not settle the complete
ontological relation between expression and determination.

Several questions remain open:

\begin{enumerate}
    \item Is expression itself a determination,
    or does it occupy a distinct ontological category?

    \item Can expression occur without being
    constitutively related to a determination?

    \item Can a determination exist without any
    corresponding expression?

    \item Under what conditions does expression
    articulate, depend upon, or contribute to
    determination?

    \item Can the relation between expression
    and determination be reciprocal?

    \item Do different domains of existence involve
    different relations between expression
    and determination?

    \item What distinguishes an ontological relation
    between expression and determination from
    a merely descriptive correspondence?
\end{enumerate}

These questions require further examination.
They cannot be resolved by assuming that expression
is simply the external appearance of determination
or that determination is merely what expression
makes available.

\subsection{Transition to Expression and Invariance}

The distinction between expression and determination
raises a further problem concerning invariance.

If an existent is associated with determinations
and expressions that may undergo change, it becomes
necessary to examine whether any aspects of expression
remain invariant across specified transformations.

However, the invariance of an expression must not
be confused with the invariance of a determination,
and neither should be treated as equivalent to
the identity of the existent.

The next section examines the relation between
expression and invariance while preserving the
distinctions among identity, determination,
expression, and persistence.

% =============================================================================
% End of Section 03
% =============================================================================